\section{Data and Methods}
\label{sec:data-method}

\subsection{Data Sources and Analytical Samples}
\label{subsec:dataset}

We study everyday moral judgments in r/AmItheAsshole (AITA), where an
original poster (OP) describes an interpersonal situation and commenters
judge the OP's conduct. We extend the analysis to r/AITAH (AITAH), which
supports a similar narrative-and-judgment format and verdict vocabulary.
AITA is the primary setting; AITAH provides a second community context in
which to examine the reception and expression of blaming and defending dissent.

We use archived public submissions and comments covering 2018--2025.
\textbf{[TODO: identify the archive source, provide its citation, and report
retrieval dates and known coverage limitations.]}
Within-community analyses use the available observations for each community.
Direct cross-community comparisons use the common period from 17 March 2021
to 23 December 2025. For each eligible verdict comment, we retain its post
identifier, text, verdict group, score, timestamp, and links to its replies.

\paragraph{Eligibility and coverage.}
Our judgment-level unit is a top-level verdict comment: a direct reply to a
submission whose extracted label belongs to a retained verdict group. The main analyses include posts
with at least $200$ retained verdict comments and exclude exact ties between
YA and NA, as defined below. Nested replies are analyzed as responses to
judgments rather than included in the post's verdict distribution.
\textbf{[TODO: specify treatment of deleted text, bots, and unavailable authors;
report counts before and after filtering by community.]}

Comment-count and score analyses use all eligible observations in the
collected dataset. This describes coverage of retained archive data, not
complete coverage of community activity. Reply depth and subtree score use a
random sample of posts because they require recursive reply-tree traversal.
Language analyses use separate annotated samples, not annotations of the full
dataset. Certainty annotation concerns top-level comments attached to sampled
posts with at least $200$ retained YA/NA verdict comments
($\mathrm{YANA\_total}\ge200$), rather than a single label for each post.
The sampling plan supplements earlier annotations toward coverage of $3\%$
of eligible high-agreement posts and $15\%$ of eligible low-agreement posts.
The certainty-annotated dataset used in this paper comprises $1{,}195$
low- and high-agreement posts. Intermediate-agreement posts present in the
broader annotation file are excluded from this dataset. Posts are the
sampling unit and comments the annotation unit; these annotations do not
cover the full collected dataset.

The $1{,}195$ posts include five exact-tie AITA posts under the stored
final-verdict metadata, which are excluded from directional comparisons.
The full-observed-period certainty analysis below further requires an observed
minority comment with a valid certainty label and includes $1{,}169$
contributing posts and $104{,}135$ minority comments. These analysis-specific
counts are distinct from the size of the annotated dataset.
\textbf{[TODO: reconcile the within-dataset breakdown: the stored metadata
give AITA 659/276 and AITAH 202/58 high/low posts, whereas the previously
reported counts are AITA 659/277 and AITAH 201/58; both total 1,195.
Document the sampling history and final-verdict metadata version. Report
completed annotation counts for the 1,195-post low/high dataset, separately
from retained analysis comments; do not use comment totals from the broader
annotation file or infer population sampling weights from nominal fractions.]}

Direct-response annotation uses three samples of $16{,}000$ parent--reply
pairs each, excluding OP participation. These are distinct from the
post-based certainty sample.
\textbf{[TODO: provide a sample-accounting table with post and comment counts,
temporal coverage, and agreement/direction strata for each analysis. Document
any analysis-specific filtering of the annotated dataset. Document
the reply-pair selection procedures and overlap among the three samples and the
number of unique pairs; confirm the reply-tree sampling design and size.
Do not describe the three samples as 48,000 unique pairs without checking.]}

\subsection{Operationalizing Minority Moral Dissent}
\label{subsec:method}

\paragraph{Verdict extraction and grouping.}
We derive each comment's \texttt{label} from the first recognized verdict
expression appearing in its \texttt{body}, not necessarily at the beginning
of the comment. Short codes are matched case-sensitively with word boundaries
(e.g., \texttt{NTA}, \texttt{NAH}, and \texttt{INFO}). This prevents lowercase
uses such as ``nah'' or ``info'' and substrings such as ``ESH'' in ``FRESH''
from being recognized as verdict codes. Full verdict phrases, such as
``not the asshole'' and ``you're the asshole'', are matched case-insensitively.
If multiple recognized verdict expressions occur, the first in textual order
determines the label.

Hypothetical forms are normalized as YWBTA to YTA and YWNBTA to NTA;
H-suffixed aliases in AITAH are mapped to their corresponding tags. After
normalization, we group labels by whether they assign fault to the OP:
$\mathrm{YA}=\{\text{YTA},\text{ESH}\}$ and
$\mathrm{NA}=\{\text{NTA},\text{NAH}\}$.
Comments labeled INFO are excluded because they request information rather
than express a verdict. Appendix~\ref{sec:appendix-tags} lists the surface
forms and their meanings.

\paragraph{Alignment and agreement.}
For a post $p$, let $N_{\mathrm{YA}}(p)$ and $N_{\mathrm{NA}}(p)$ count
retained top-level verdicts in each group. We define the post's final
\emph{YA share}, denoted $\rho(p)$, as
\[
\rho(p)=\frac{N_{\mathrm{YA}}(p)}{N_{\mathrm{YA}}(p)+N_{\mathrm{NA}}(p)}.
\]
Throughout, ``YA share'' refers to this proportion of YA judgments, and
$\rho(p)$ denotes its final observed value for post $p$. Time-dependent
proportions are explicitly described as running or window-specific shares;
minority share instead counts the non-dominant verdict group.
Posts with $\rho(p)=0.5$ are excluded. The dominant verdict is
$D(p)=\mathrm{YA}$ when $\rho(p)>0.5$ and $D(p)=\mathrm{NA}$ otherwise.
A comment is majority-aligned when its verdict matches $D(p)$ and
minority-aligned otherwise. High-agreement posts satisfy $\rho(p)\le0.20$
or $\rho(p)\ge0.80$; low-agreement posts satisfy
$0.40\le\rho(p)\le0.60$, excluding exact ties. Intermediate bands are omitted
from this comparison. Low-agreement posts provide a baseline comparison,
while the primary dissent analyses focus on high-agreement posts.

\paragraph{Direction of dissent.}
Blaming dissent is a minority YA judgment in an NA-dominant post; defending
dissent is a minority NA judgment in a YA-dominant post. Each direction is
compared with majority comments in the corresponding post type. The
between-direction comparison uses separate NA-dominant and YA-dominant post
cohorts. All alignment and agreement labels use the final observed verdict
distribution, including for annotated posts, rather than a measure of
commenters' contemporaneous perceptions.
Table~\ref{tab:dimensions} summarizes the definitions and measures.

\begin{table*}[t]
  \centering
  \small
  \begin{tabular}{@{}p{0.16\textwidth}p{0.20\textwidth}p{0.58\textwidth}@{}}
    \toprule
    Variable & Levels / range & Definition \\
    \midrule
    \rowcolor{gray!15}
    \multicolumn{3}{@{}l}{\emph{Analytical dimensions}} \\
    \addlinespace[2pt]
    Alignment & majority / minority & Whether a verdict comment matches the post's final observed dominant verdict $D(p)$. \\
    Agreement & {high / low} & High: $\rho(p)\le0.2$ or $\ge0.8$; low: $0.4\le\rho(p)\le0.6$. Exact ties ($\rho(p)=0.5$) are excluded. \\
    Dissent direction & blaming / defending & Minority YA in NA-dominant posts / minority NA in YA-dominant posts. \\
    Community & AITA / AITAH & Source subreddit. \\
    \midrule
    \rowcolor{gray!15}
    \multicolumn{3}{@{}l}{\emph{Measures}} \\
    \addlinespace[2pt]
    Comment score & integer & Last available net score (upvotes minus downvotes) of the verdict comment in the dataset. \\
    Direct-reply count & integer $\ge0$ & Number of recorded direct replies to the verdict comment. \\
    Reply depth & integer $\ge0$ & Maximum recorded reply depth below the verdict comment; $0$ for no replies. \\
    Subtree score & integer & Sum of recorded scores of the verdict comment and all replies in its observed subtree. \\
    Direct-response type & four categories & Parent--reply annotation: support, challenge, information-seeking, or other. The child is classified relative to its immediate parent. \\
    Expressed certainty & ordinal, $1$--$4$ & Annotated certainty of a verdict comment, from unresolved judgment ($1$) to a confidently asserted judgment ($4$); measured only in the certainty-annotated sample and its specified subsets. \\
    \bottomrule
  \end{tabular}
  \caption{Analytical dimensions and measures. $\rho(p)$ denotes the final
  YA share of post $p$.
  Minority status uses the final observed verdict distribution.
  Direct-response annotation excludes OP participation; its sample is
  separate from the certainty sample. Sampling and aggregation procedures are
  described in the text.}
  \label{tab:dimensions}
\end{table*}


\subsection{Reception Measures and Response Annotation (RQ1)}
\label{subsec:method-reception}

\paragraph{Evaluation and conversational engagement.}
Comment score captures the recorded net voting outcome. Direct-reply count
measures replies to the verdict comment; reply depth measures the maximum
number of nested reply levels beneath it. Subtree score sums scores over the
verdict comment and its reply subtree. We calculate this aggregate separately
from the initiating comment's own score. Analyses distinguish reply incidence from the amount
or depth of engagement, with the relevant reply-count or depth threshold
specified for each result. Majority--minority comparisons across agreement
levels establish the baseline; the primary comparison concerns blaming and
defending dissent and their respective majority baselines.
We use the final score available for each comment in the dataset, rather than
measuring scores at a fixed interval after comment creation. Reply measures
use all recorded replies linked to eligible comments, without imposing an
additional elapsed-time cutoff for these aggregate reception analyses.
These outcomes reflect the available archive records, not necessarily the
eventual lifetime scores or complete reply histories. Comment-level means
and proportions are pooled over eligible comments in the relevant analytical
sample, giving each comment equal weight rather than first averaging within
posts. Posts with more eligible comments therefore contribute more to these
summaries. Score and direct-reply summaries use all eligible comments in the
collected dataset; the separately annotated response sample is described below.

\paragraph{Direct-response annotation.}
\label{subsec:method-rq2}
We classify a direct child reply's response to its immediate parent verdict
comment, using the original post as context. The four categories are
agreement/support, disagreement/challenge, clarification/information-seeking,
and other. Labels concern communicative function rather than sentiment:
a polite rebuttal can be a challenge, while qualified agreement can remain
support. The prompt does not supply comment scores, minority/majority labels,
or agreement-level metadata. Its category definitions and rules for mixed
responses appear in Appendix~\ref{sec:appendix-rq2-annotation}.

The annotated response analysis excludes OP participation, distinguishing
responses by other commenters from engagement by the person being judged.
\textbf{[TODO: specify whether the OP exclusion applies to parent authors,
reply authors, or both, and how unavailable author identities are handled.]
}
We compare response categories across parent alignment, dissent direction,
and community. Challenge proportions use labeled direct replies as the
denominator; reply incidence uses eligible parent verdict comments, including
those without replies. These measures are computed separately.
\textbf{[TODO: report the sampling strata and any weights, model/version and
decoding settings, human-validation sample and protocol, and category-level
performance. Distinguish prompt-development cases from held-out validation,
and missing labels from the other category.]}

\subsection{Minority Prevalence Over Time (RQ2)}
\label{subsec:method-temporal}

We examine the prevalence of blaming and defending dissent as discussions
unfold, using eligible top-level verdict comments rather than restricting
participation analyses to the certainty-annotated sample. High-agreement
posts are the primary focus; low-agreement posts provide a comparison.
Time is measured from submission to comment arrival.

\paragraph{Cumulative and window-specific minority share.}
We distinguish window-specific minority share---the fraction of verdict comments arriving
in a given window that express the relevant minority judgment---from
cumulative share, which includes comments up to that time. Blaming share is
computed within NA-dominant posts and defending share within YA-dominant
posts, retaining the final-distribution labels throughout. Cumulative and
window-specific shares address different aspects of participation and are
reported separately. The displayed time range is specified for each figure.
\textbf{[TODO: verify the cohort, time bins, comment-versus-post weighting,
and uncertainty procedure for each share figure. Report contributing post
and comment counts and the treatment of empty windows and truncated follow-up.]}
Changes describe the composition of observed judgments, not a within-person
panel: contributing posts and commenters can differ between windows.

\subsection{Expressed Linguistic Certainty (Supplementary)}
\label{subsec:method-certainty-contrasts}

Separately from temporal participation, we characterize how confidently the
two directions of dissent are expressed in the annotated low/high sample.
The ordinal scale runs from unresolved judgment ($1$) to a confidently
asserted single judgment ($4$); its prompt and definitions appear in
Appendix~\ref{sec:appendix-annotation}. This measures expressed language,
not latent attitude certainty. Verdict extraction still follows the
first-match rule, even for comments annotated as unresolved.

We use the full observed period without an upper elapsed-time cutoff,
retaining YA/NA minority comments with valid certainty labels in
$\{1,2,3,4\}$. Within each post, we calculate the proportions at each level
and then average them with equal post weight. Posts without a valid minority
certainty observation do not contribute and are not assigned zero.
We summarize directional differences using the certainty-$4$ proportion,
with $95\%$ percentile intervals from $10{,}000$ whole-post bootstrap
replicates, independently resampling the two directional cohorts.
Full-period distributions, mean-score contrasts, and annotation-version
and first-18-hour sensitivity checks are reported in
Appendices~\ref{sec:appendix-certainty-comparison}
and~\ref{sec:appendix-certainty-window}. These are comparisons between
post cohorts, not estimates of certainty changing over time. Archive
follow-up may differ across posts and need not cover complete lifetimes.

\paragraph{Descriptive scope and uncertainty.}
Means over all eligible archive observations are descriptive quantities for
that retained dataset. They are distinguished from estimates based on
sampled posts or annotated replies. For the reported post-bootstrap intervals,
posts are the resampling unit, retaining within-post dependence among comments.
\textbf{[TODO: document the uncertainty procedures for the remaining
analyses and any additional effect-size or interaction comparisons. Do not imply that all
figures use bootstrap intervals or that an annotated sample is representative
without documenting its selection.]}

\subsection{Cross-Community Extension (RQ3)}
\label{subsec:method-community}

We apply the same verdict definitions and outcome measures in each community
and compare reception measures and the temporal prevalence of dissent.
Direct comparisons use the common calendar period specified in
Section~\ref{subsec:dataset}, with comment-arrival windows reported separately.
Community membership is an observed grouping variable, not an experimentally
assigned condition; no causal effect of a verdict rule is estimated.

To characterize content overlap, we use post embeddings from
\texttt{BAAI/bge-base-en-v1.5} \citep{bge_embedding} and
\texttt{intfloat/e5-base-v2} \citep{wang2022text}, examining projections,
community-prediction performance, distributional distances, and topic
composition (Appendix~\ref{sec:appendix-content}). We also describe shared
users and authors using intersection counts and Jaccard overlap, including
activity- and tenure-restricted cohorts (Appendix~\ref{sec:appendix-overlap}).
We also summarize the distributions of elapsed time to comment arrival
(Appendix~\ref{sec:appendix-timing}).
\textbf{[TODO: document embedding preprocessing, classifier evaluation,
permutation settings, topic-model specification, and user-cohort definitions
in the corresponding appendices.]}

\paragraph{Cross-posted stories.}
Cross-posted stories are reserved for illustrative examples rather than a
principal quantitative test. Candidate matches are identified from identical
title-and-body text after Unicode NFKC normalization, case folding, and
whitespace normalization; empty, deleted, and removed bodies are excluded.
We identify $4{,}150$ distinct stories, defined by unique normalized content,
that appear in both AITA and AITAH. Of these, $3{,}913$ have exactly one
submission in each community and form one-to-one pairs. The remaining $237$
appear multiple times in at least one community, allowing multiple possible
submission pairings for the same story.

Among the cross-community submission pairs, $1{,}871$ have an observed YA
share $\rho(p)$ in both communities. Submissions with no retained verdicts
($\mathrm{YANA\_total}=0$) have an undefined YA share; these missing values
are not imputed. Requiring at least $20$ retained YA/NA verdicts in each
submission ($\mathrm{YANA\_total}\ge20$ in both communities) leaves $107$
pairs. Here, $\mathrm{YANA\_total}=N_{\mathrm{YA}}(p)+N_{\mathrm{NA}}(p)$.
This case-selection threshold is separate from the $200$-verdict threshold
used in the main analyses.

\textbf{[TODO: specify how the 237 multiply submitted stories enter the
candidate pair set, and select and verify illustrative cases from the
107 retained pairs. Report case-selection criteria and verdict counts;
examples must not be presented as representative evidence of a community
effect.]}
